\chapter{Ingestor agéntico: LLM-as-compiler offline y drift-as-signal}
\label{sec:paper-ingestor}

\section{Problema}

La ingesta productiva de logs tiene dos propiedades en tensión:
debe ser \emph{determinística} (dos corridas del mismo batch
producen la misma salida, byte-a-byte) y debe ser \emph{adaptable}
(formatos nuevos aparecen, y los existentes evolucionan aguas
arriba sin aviso). La práctica dominante resuelve la primera y
abandona la segunda: un parser hardcodeado por formato, release del
daemon por cada cambio, parser genérico heurístico como red de
seguridad con pérdida de fidelidad silenciosa.

GraphHunter resuelve las dos simultáneamente separando el sistema
en tres carriles con perfiles distintos de determinismo, latencia y
quién los opera. El presente capítulo describe el diseño, la
matemática del contrato, y qué asume sobre el humano en el loop.

\section{Decisión: tres carriles}

\begin{itemize}
    \item \textbf{Fast lane}. Un compilador determinístico
    $\text{compile}: \text{FieldConfig} \to \text{VRL}$ emite
    programas textuales con identidad estable
    $\text{hash}(f) = \text{sha256}(\text{compile}(f))$. El runtime
    ejecuta VRL por evento. Cero LLM, cero red, cero
    nondeterminismo. El hash se estampa en cada evento canonical
    (\code{provenance.mapping_hash}), cerrando la trazabilidad
    dato$\leftrightarrow$lógica.
    \item \textbf{Slow lane}. Un LLM actúa como \emph{compilador
    asistente} offline: dado una muestra del formato y un target
    canonical (OCSF v1.4), propone un \code{FieldConfig}. El output
    del LLM son datos, no código. El código lo produce el compilador
    VRL determinístico a partir de esos datos.
    \item \textbf{Learning lane}. El fast lane emite métricas
    estructurales baratas por batch (histogramas, cardinalidades,
    ratios DLQ). Un monitor de drift compara contra el perfil
    esperado del mapping activo; cuando supera umbral, abre
    re-review en el slow lane sin detener la ingesta.
\end{itemize}

El contrato entre carriles es un par
$(\text{FieldConfig}, \text{OCSFMapping})$: ambos datos, ambos
versionados. Ningún carril llama a otro por función en tiempo real.
El acoplamiento es por artefacto.

\section{LLM como compilador, no como traductor}

La literatura de ingesta con LLM se suele deslizar a una variante
riesgosa: el LLM procesa eventos en caliente. Esto falla de
manera predecible por tres razones:

\begin{enumerate}
    \item \textbf{Latencia}. Un LLM remoto agrega cientos de ms por
    request. Un LLM local agrega decenas de ms. Para un stream de
    miles de eventos por segundo, ninguno alcanza.
    \item \textbf{Costo}. Por evento ingerido pagás un token-count.
    A escala SIEM (millones de eventos/día), este costo es órdenes
    de magnitud mayor que el compute determinístico.
    \item \textbf{Superficie de falla}. Un LLM que parsea un batch
    mal aplica el error a todos sus eventos. Sin gate humano, ese
    error entra al grafo y contamina todo hunt posterior.
\end{enumerate}

La decisión de diseño: el LLM no ve eventos productivos. Ve
muestras explícitas que el operador o el monitor de drift le
entregan. Lo que produce son \emph{datos que describen una
transformación} (el FieldConfig), no la transformación ejecutada.
La ejecución la hace un compilador determinístico que no es un
LLM.

Esto es análogo a cómo un compilador optimizador corre offline y
produce código que un runtime simple ejecuta en caliente. Aquí el
compilador offline es un LLM asistido por humano; el runtime
caliente es el VRL interpreter.

\section{Drift como señal, no como error}

Un cambio aguas arriba (nuevo campo, separador modificado, semántica
movida) es, desde la perspectiva del ingestor, un cambio en la
distribución empírica de la entrada. Si lo tratamos como
\emph{error}, las opciones son dos: parar la ingesta (inaceptable en
producción), o ingerir basura (inaceptable para análisis).

Lo tratamos como \emph{señal}. Operacionalmente:

\begin{itemize}
    \item Cada batch $B$ genera un vector empírico
    $\hat{P}(B) \in \mathbb{R}^k$ de estadísticas (histograma por
    tipo, cardinalidad por campo, ratio DLQ, distribución de
    timestamps).
    \item El mapping $m$ en uso tiene un perfil esperado $P_m$
    calculado sobre su ventana de calibración.
    \item Una prueba Kolmogorov-Smirnov por campo, agregada por
    max-divergencia, produce $d(B, P_m)$.
    \item Si $d(B, P_m) > \tau$, se abre una \emph{propuesta de
    re-review} — entrada en la cola del slow lane con el batch
    anómalo adjunto. El fast lane sigue ingestando $B$ con el
    mapping $m$.
    \item El humano, con ayuda del LLM, decide si el mapping debe
    evolucionar. En el caso afirmativo, publish\_mapping emite un
    nuevo $m'$ con hash distinto; el fast lane recoge $m'$ en el
    siguiente batch.
\end{itemize}

Esto invierte el esquema tradicional. En lugar de parar para
aprender, aprendemos mientras seguimos ingestando. El costo es que
durante el período de drift detectado pero no aprobado, el grafo
recibe eventos con mapping potencialmente degradado; la mitigación
es que $\tau$ se configura bajo y la ventana de revisión es
minutos-horas, no días.

\section{Propiedades formales}

Sea $f \in \text{FieldConfig}$ y $e$ un evento crudo. Escribamos
$\text{run}(f, e) = \text{vrl\_run}(\text{compile}(f), e)$ para la
salida canonical del fast lane.

\begin{description}
    \item[Determinismo.] Para todo $f$, $e$:
    $\text{run}(f, e)$ depende sólo de $f$ y $e$. La misma entrada
    produce idéntico output en corridas distintas y máquinas
    distintas.
    \item[Trazabilidad.] Cada evento canonical lleva
    $\text{hash}(f)$ en su campo \code{provenance.mapping\_hash}. A
    partir del hash se recupera $f$ desde la mapping library; a
    partir de $f$ se reconstruye $\text{compile}(f)$; ejecutar
    sobre el evento crudo original reproduce el canonical.
    \item[Gate humano.] Toda transición
    $m \to m'$ en la mapping library es mediada por aprobación
    humana (\code{agentic_review_approve}). No hay camino
    automático LLM$\to$productivo.
    \item[No-block.] El learning lane no detiene el fast lane. Si
    el LLM falla, se cae la red, o el humano tarda, el fast lane
    sigue ejecutando el $m$ vigente.
\end{description}

Las cuatro propiedades son arquitectónicas, no validadas
empíricamente en este reporte. La evaluación cuantitativa
(latencia del fast lane comparada con parsers hardcoded, accuracy
del LLM en inferir FieldConfig, tasa de verdaderos/falsos
positivos del drift monitor) es trabajo futuro y se reporta en
capítulos separados cuando haya números.

\section{Trabajo relacionado}

El uso del LLM como compilador asistido tiene precedentes en
\emph{code synthesis} (Copilot, Codex) y \emph{schema matching} con
transformers. La diferencia aquí es que el compilador humano-LLM
produce \emph{datos tipados} (\code{FieldConfig}), no código de
programa; y la ejecución en caliente es determinística sobre esos
datos.

Drift detection sobre streams tiene tradición en machine learning
de producción (ADWIN, Page-Hinkley). El aporte de GraphHunter es
aplicar estas métricas no a predicciones sino a \emph{forma
estructural} (histogramas de tipo, cardinalidades) como proxy
barato de cambio de schema.

OCSF como target canonical es adopción de un estándar externo; el
aporte de GraphHunter es usarlo como \emph{lingua franca} entre
mappings heterogéneos y el motor de grafo.

\section{Limitaciones y trabajo futuro}

\begin{itemize}
    \item El fast lane usa un intérprete VRL embebido; la sustitución
    por la crate \code{vrl} upstream es follow-up de performance.
    \item No hay evaluación cuantitativa aún. Los números de
    latencia, accuracy LLM, y drift precision/recall se reportan
    cuando el sistema lleve tiempo suficiente en producción.
    \item El umbral $\tau$ de drift es global; idealmente sería
    por-campo y auto-ajustable a partir de señal de false-positive.
    \item El gate humano es el ancla de seguridad; también es el
    cuello de botella. Optimizar la UX de review agéntica es
    trabajo de producto, no de arquitectura.
\end{itemize}
